\documentclass[10pt,a4paper]{amsart}
\usepackage[margin=19mm]{geometry}
\usepackage{amsmath,amssymb,mathtools}
\usepackage[scr=rsfs,cal=euler]{mathalfa}
\usepackage{xfrac}
\usepackage[unicode,hidelinks,bookmarks=false]{hyperref}
\newcommand{\ie}{i.e.\ }
\newcommand{\cf}{cf.\ }
\newcommand{\resp}{respectively\ }
\newcommand{\Subsection}{\subsection}
\providecommand{\nobreakdash}{\nobreak\hbox{-}}
\setlength{\emergencystretch}{2em}
\multlinegap0pt
\usepackage{dsfont}
\usepackage{braket}
\usepackage{amssymb}
\usepackage{xcolor}
\usepackage{enumerate}
\usepackage{tikz}
\usepackage{mathtools}
\usepackage{bm}
\usepackage[normalem]{ulem}
\newtheorem{prop}{Proposition}
\newcommand{\defn}{:=}
\title{Distinct Types of Parent Hamiltonians for Quantum States:\\ Insights from the $W$ State as a Quantum Many-Body Scar}
\author{Lei Gioia, Sanjay Moudgalya, Olexei I. Motrunich}
\date{Source: arXiv:2510.24713v3 (CC BY 4.0)}
\begin{document}
\maketitle

\noindent\emph{Source and licence.} Distinct Types of Parent Hamiltonians for Quantum States:\\ Insights from the $W$ State as a Quantum Many-Body Scar. Version arXiv:2510.24713v3 (CC BY 4.0), distributed by arXiv under the Creative Commons Attribution 4.0 International licence, http://creativecommons.org/licenses/by/4.0/, as recorded in the arXiv licence metadata for this exact version and shown on the abstract page https://arxiv.org/abs/2510.24713v3. Full licence notice: \texttt{licenses/arxiv-2510.24713-CC-BY-4.0.txt}.\par\smallskip

\noindent\textit{Editorial scope.} Counted unit: the article's prop prop:W2 (source0.tex lines 1189--1192). The article's complete original proof of it is delivered with it (source0.tex lines 3019--3117, one \texttt{proof} environment of 99 lines, which the article places away from the statement). No mathematics is written, repaired or re-derived: the statement and the proof are the article's own lines, in the article's own notation, and no retained line is substituted or re-flowed.  Retained uncounted as the source-local context the proof needs: lines 739--742; lines 949--952.  Nothing was omitted from any retained span, so the package carries no omission record.  No retained line contains a citation, so the wrapper carries no bibliography and no bibliography entry.  No retained line contains a figure environment, an image-inclusion command or drawing code, so no image is delivered and the package declares no asset dependency.  Editorial formatting only, in addition: an AMS \texttt{amsart} preamble that restates the article's own declarations and macros used by the retained lines, namely the environments \texttt{prop} on separate counters, together with the standard packages those lines need.  The retained environments print in this document as Proposition 1 in order of appearance, whereas the article prints the same items with the numbering of its own full document.  The complete native source of the article as distributed in the arXiv e-print for this exact version is bundled unchanged as \texttt{sources/187982/source0.tex}.
\begin{align}
    H =  \Omega \mathds{1} + \omega \hat{N}_{\rm tot} + t H_{\rm ImHop} + \sum_{X,\,|X|\leq R_{\text{max}}} h_X \quad,
    \label{eq:HW}
\end{align}

\begin{equation}
    \ket{W^2} \defn \frac{1}{\sqrt{\binom{N}{2}}} \sum_{j < k} s_j^\dagger s_k^\dagger \ket{\bar{0}}\quad,
    \label{eq:W2}
\end{equation}

\begin{prop}
\label{prop:W2}
    $\ket{W^2}$, defined in Eq.~\eqref{eq:W2}, is an asymptotic scar with the energy variance upper-bounded by $O(N^{-1})$.
\end{prop}

\begin{proof}
    Our arguments roughly follow the proof of Prop.~\ref{prop:W2} in the main text of $\ket{W^2}$ as an asymptotic scar, with new technical steps needed for $p \geq 3$.
    %
    We start with writing of the parent Hamiltonian given by Eq.~(\ref{eq:HW}).
    %
    Consider a given $h_X$ with bounded range and bounded norm.
    %
    The Schmidt decomposition of $\ket{W^p}$ for a bipartition into regions $X$ and its compliment $X^c$ is given by
    \begin{equation}
        \ket{W^p} = \sum_{l=0}^{p}  f_{l}^{X,N} \ket{W^l}_X\otimes \ket{W^{p-l}}_{X^c} ~, \qquad
        f_{l}^{X,N} \defn \sqrt{\frac{\binom{|X|}{l}\binom{N-|X|}{p-l}}{\binom{N}{p}}} \approx \sqrt{\frac{\binom{|X|}{l} p!}{(p-l)!}}\, \frac{1}{\sqrt{N^l}} ~,
        \label{eq:flapprox}
    \end{equation}
    where we have assumed for simplicity $|X| \geq p$ (otherwise, the sum over $l$ terminates at $|X|$).
    %
    In the second equation, we have also shown the behavior of the amplitudes $f_{l}^{X,N}$ for large $N$, assuming $|X|$ and $p$ are fixed ($l \leq \text{min}\{|X|,p\}$).
    %

    %
    Upon application of $h_X$, recall that $h_X\ket{\bar{0}}_X=h_X\ket{W}_X=0$, such that
    %
    \begin{equation}
    h_X \ket{W^p} = \sum_{l=2}^{p} f_{l}^{X,N} h_X \ket{W^l}_X\otimes \ket{W^{p-l}}_{X^c}\,\;\text{and}\;\;\langle h_X \rangle_{W^p} = \sum_{l=2}^p (f_{l}^{X,N})^2 \bra{W^l}_X h_X \ket{W^l}_X ~,
    \label{eq:hXWp}
    \end{equation}
    where $\langle...\rangle_{W^p}:=\bra{W^p}...\ket{W^p}$.
    %
    This is upper-bounded by
    \begin{equation}
    |\langle h_X \rangle_{W^p}| \leq \|h_X\| \sum_{l=2}^p (f_{l}^{X,N})^2 \approx \|h_X\| (f_{2}^{X,N})^2
    \approx \|h_X\| \frac{|X|(|X|-1) p(p-1)}{2 N^2}
    \end{equation}
    for large $N$.
    %
    Hence,
    \begin{equation}
    \big|\big\langle \sum_X h_X \big\rangle_{W^p}\big| \leq \frac{C}{N} ~.
    \label{eq:h_XWpexp}
    \end{equation}
    %
    Thus, the energy of $\ket{W^p}$ (as a trial state) can be shifted from $p \omega$ by an amount that decreases as $1/N$ (where we have used that $H_{\text{ImHop}} \ket{W^p} = 0$ and $N_{\text{tot}} \ket{W^p} = p \ket{W^p}$).
    %

    %
    Consider now the calculation of the variance.
    %
    In the case of $p=2$ in the previous section, we had $h_Y h_X \ket{W^p} = 0$ if $X \cap Y = \emptyset$.
    %
    This holds also for $p = 3$ but is no longer true for $p \geq 4$.
    %
    Nevertheless, we can controllably analyze such contribution as follows.
    %
    Assuming $X \cap Y = \emptyset$ we apply $h_Y$ to the left equation in Eq.~(\ref{eq:hXWp}) to obtain
    \begin{equation}
        h_Y h_X\ket{W^p} = \sum_{l,m=2}^{p} f_{l}^{X,N}f_{m}^{Y,N-|X|} h_X\ket{W^l}_X \otimes h_Y\ket{W^m}_Y \otimes \ket{W^{p-l-m}}_{(X\cup Y)^c} ~,
    \end{equation}
    where we have successively Schmidt decomposed the state, first into $X$ and $X^c$, and then we have further partitioned $X^c$ into $Y$ and $(X \cup Y)^c$.
    %
    It follows that
    \begin{equation}
        \langle h_Y h_X \rangle_{W^p} = \sum_{l,m=2}^p f^{X\cup Y,N}_{l+m}f_{l}^{X,N}f^{Y,N-|X|}_{m} \,
        \bra{W^{l+m}}_{X\cup Y}h_Y h_X \ket{W^l}_X \otimes \ket{W^m}_Y ~.
    \end{equation}
    The above equation for $X \cap Y = \emptyset$ is upper-bounded by
    \begin{equation}
        |\langle h_Y h_X \rangle_{W^p}| \leq \|h_X\| \|h_Y\| \sum_{l,m=2}^p f^{X\cup Y,N}_{l+m}f_{l}^{X,N}f^{Y,N-|X|}_{m} 
        = O(N^{-4}) ~,
        \label{eq:hyhxnooverlap}
    \end{equation}
    where the leading scaling with $N$ comes from the $l = m = 2$ term.

    %
    On the other hand, when $X \cap Y \neq \emptyset$, we can write, e.g., 
    \begin{equation}
        h_Y h_X \ket{W^p} = \sum_{l=2}^p f_l^{X\cup Y,N} h_Y h_X \ket{W^l}_{X\cup Y} \otimes \ket{W^{p-l}}_{(X\cup Y)^c} ~,
    \end{equation}
    such that
    \begin{equation}
        \langle h_Y h_X\rangle_{W^p}=\sum_{l=2}^{p} \left(f_{l}^{X\cup Y,N}\right)^2\bra{W^l}_{X\cup Y}  h_Y h_X\ket{W^l}_{X\cup Y}.
    \end{equation}
    Such an expectation value is upper-bounded by
    \begin{equation}
        |\langle h_Y h_X\rangle_{W^p}| \leq \sum_{l=2}^{p} \left(f_{l}^{X\cup Y,N}\right)^2 \|h_Y h_X\|
        \approx \left(f_{2}^{X\cup Y,N}\right)^2  \|h_Y h_X\| = O(N^{-2}) ~.
        \label{eq:hyhxoverlap}
    \end{equation}
    Putting Eqs.~(\ref{eq:hyhxnooverlap}) and (\ref{eq:hyhxoverlap}) together, we can upper-bound
    \begin{equation}
        \big|\big\langle \big(\sum_X h_X \big)^2\big\rangle_{W^p}\big| \leq O(N^{-1}) ~,
    \end{equation}
    since there are $O(N^2)$ contributions from $X \cap Y = \emptyset$, each bounded by $O(N^{-4})$, and $O(N)$ contribution from $X \cap Y \neq \emptyset$, each bounded by $O(N^{-2})$.
    The energy variance of $\ket{W^p}$, calculated using the above equation and Eq.~(\ref{eq:h_XWpexp}), is
    \begin{align}
        \big\langle \big(\sum_X h_X \big)^2\big\rangle_{W^p}-\big\langle \sum_X h_X \big\rangle_{W^p}^2 \leq O(N^{-1}) \quad.
    \end{align}
    %
    This gives a lifetime of $\ket{W^p}$ to also be $\sim \sqrt{N}$, where $p \ll N$ to have controlled approximations such as in Eq.~(\ref{eq:flapprox}).
    %
\end{proof}

\end{document}
